\documentclass[10pt,a4paper]{amsart}
\usepackage[margin=19mm]{geometry}
\usepackage{amsmath,amssymb,mathtools}
\usepackage[scr=rsfs,cal=euler]{mathalfa}
\usepackage{xfrac}
\usepackage[unicode,hidelinks,bookmarks=false]{hyperref}
\newcommand{\ie}{i.e.\ }
\newcommand{\cf}{cf.\ }
\newcommand{\resp}{respectively\ }
\newcommand{\Subsection}{\subsection}
\providecommand{\nobreakdash}{\nobreak\hbox{-}}
\setlength{\emergencystretch}{2em}
\multlinegap0pt
\usepackage{bm}
\usepackage{bbm}
\usepackage{dsfont}
\usepackage{xspace}
\usepackage{cancel}
\usepackage{mathtools}
\usepackage{amsthm}
\makeatletter
\@ifundefined{theorem}{\newtheorem{theorem}{Theorem}}{}
\@ifundefined{corollary}{\newtheorem{corollary}[theorem]{Corollary}}{}
\@ifundefined{lemma}{\newtheorem{lemma}[theorem]{Lemma}}{}
\@ifundefined{proposition}{\newtheorem{proposition}[theorem]{Proposition}}{}
\makeatother
\newcommand\la{\lambda}
\newcommand{\D}{\mathbb{D}}
\newcommand{\T}{\mathbb{T}}
\title[Boundary behavior of functions in the Schur-Agler class of t]{Boundary behavior of functions in the Schur-Agler class of the polydisc}
\author{Jim Agler, Connor Evans, Zinaida Lykova, N. J. Young}
\date{Source: arXiv:2508.13742v2 (CC BY 4.0)}
\begin{document}
\maketitle

\noindent\emph{Source and licence.} Boundary behavior of functions in the Schur-Agler class of the polydisc. Version arXiv:2508.13742v2 (CC BY 4.0), distributed under the Creative Commons Attribution 4.0 International licence, as recorded in the arXiv licence metadata for this exact version and shown on the abstract page https://arxiv.org/abs/2508.13742v2. Full rights notice: licenses/arxiv-2508.13742-CC-BY-4.0.txt.\par\smallskip

\noindent\textit{Editorial scope.} Counted unit: the Proposition whose source label is \texttt{prop5.8cara} in the article (source0.tex lines 1164--1167, source label \texttt{prop5.8cara}), delivered together with the article's own complete original proof of it (source0.tex lines 1168--1366, one \texttt{proof} environment of 199 lines). No mathematics is written, repaired or re-derived: statement and proof are the article's own text line for line. Required context retained uncounted: modeleqn (lines 191--193); whenlambdametmu (lines 545--549); prop4.2generalized (lines 618--630); theorem4.9generalized (lines 655--662); phi-omega (lines 699--704); phi-nt-omega (lines 763--766); coro5.7cara (lines 1105--1121); ujnotxj (lines 1127--1134). Every definition, display and label of those lines is retained unchanged. Omitted from this package and not redistributed: 1--190, 194--544, 550--617, 631--654, 663--698, 705--762, 767--1104, 1122--1126, 1135--1163, 1367--3340. Nothing inside a retained excerpt is omitted or altered: every retained body is delivered exactly as the article's lines are written, so no omission receipt of any kind accompanies this package. Citations: the retained excerpts carry no citation command at all, so no bibliography entry of the article is transcribed and no reference is redistributed. Reference adaptation: none was needed and none was made; every reference inside the delivered unit resolves inside it, so no unresolved reference or citation is produced. Printed slips kept literal: no printed word, symbol, hypothesis or number of the counted statement or of its proof was altered. Layout and class: the article's own class is \texttt{amsart} with packages including amsmath, amssymb, amsfonts, dsfont, color, version; the wrapper is the lane's pinned \texttt{amsart} editorial wrapper at 10pt on A4, which restates the theorem-like environments the retained text needs and adds the article's own macro definitions that the retained text needs (\texttt{\textbackslash D}, \texttt{\textbackslash T}, \texttt{\textbackslash la}), with extra wrapper packages (bm, bbm, dsfont, xspace, cancel, mathtools). The delivered numbering is the unit's own. Assets: no retained span contains a figure environment, a graphics-inclusion command, a PDF-page inclusion, a table or a TikZ picture; no image file is copied into this package, the package declares no asset dependency and no image file is redistributed with it. Licence: version arXiv:2508.13742v2 of this article is distributed under the Creative Commons Attribution 4.0 International licence, as recorded in the arXiv licence metadata for this exact version and shown on the abstract page https://arxiv.org/abs/2508.13742v2. A full rights notice is delivered at licenses/arxiv-2508.13742-CC-BY-4.0.txt. Characters: every retained excerpt is pure ASCII, so the delivered engine, XeLaTeX with its default Latin Modern text font, reports no missing character.
\begin{equation}\label{modeleqn}
1-\overline{\varphi(\mu)}\varphi(\lambda) = \Big\langle{\big(1_{\mathcal{L}}-\mu_{P}^{*}\lambda_{P}\big)v({\lambda}),v({\mu})} \Big\rangle_{\mathcal{L}},
\end{equation}

\begin{align}
1-\lvert\varphi(\lambda)\rvert^{2} &= \Big\langle{\big(1_{\mathcal{M}}-\sum_{j=1}^{d}\lvert \lambda_{j}\rvert^{2}P_{j}\big)u({\lambda}),u({\lambda})} \Big\rangle_{\mathcal{M}}\nonumber\\
    &=\sum_{j=1}^{d}\Big\langle (1-\lvert \lambda_{j}\rvert^{2})u_{j}({\lambda}),u_{j}({\lambda}) \Big\rangle_{\mathcal{M}_{j}}\nonumber\\
    &=\sum_{j=1}^{d}(1-\lvert \lambda_{j}\rvert^{2})\|u_{j}(\lambda)\|^{2}_{\mathcal{M}_{j}}.\label{whenlambdametmu}
\end{align}

\begin{proposition}\label{prop4.2generalized}
Let $\tau\in\partial\mathbb{D}^{d}$ be a $B$-point for $\varphi\in\mathcal{SA}_{d}$, let $(\mathcal{M},u)$ be a model of $\varphi$ and let $Y_{\tau}$ be the cluster set of the  model $(\mathcal{M},u)$ at $\tau$.
\begin{enumerate}
\item If
   $$\liminf_{\substack {\lambda{\overset{nt}\rightarrow}\tau \\ \lambda \in \mathbb{D}^{d}}} \dfrac{1-\lvert \varphi(\lambda) \rvert}{1-\|\lambda\|_{\infty}}=\alpha,$$
then there exists an $x\in Y_{\tau}$ such that $\| x\|^{2}_{\mathcal{M}}\leq \alpha$.
\item If $x\in Y_{\tau}$ and $\lvert\tau_{j}\rvert< 1$ for any $j=1,\hdots,d$, then $x_{j}=0$.
\item There exists $\omega\in\mathbb{T}$ such that, for all $x\in Y_{\tau}$ and $\lambda\in\mathbb{D}^{d}$,
    \begin{equation}\label{omegaequationintheoremstatement}
   1-\overline{\omega}\varphi(\lambda)=\sum_{\lvert \tau_{j} \rvert = 1}\Big(1-\overline{\tau}_{j}\lambda_{j}\Big)\langle u_{j}(\lambda),x_{j}\rangle_{\mathcal{M}_{j}}.
    \end{equation}
\end{enumerate} 
\end{proposition}

\begin{theorem}\label{theorem4.9generalized}
Let $\varphi\in\mathcal{SA}_{d}$, $\tau\in\partial\mathbb{D}^{d}$. If there is a sequence $\{\lambda^{(n)}\}$ in $\mathbb{D}^{d}$ converging to $\tau$ such that
$$\lim_{n\rightarrow \infty}\dfrac{1-\lvert \varphi(\lambda^{(n)})\rvert}{1-\|\lambda^{(n)}\|_{\infty}}=\alpha<\infty,$$
then there exists $\omega\in\mathbb{T}$ such that, for all $\lambda\in\mathbb{D}^{d}$, 
\begin{equation}\label{alphamaxequtionforgeneralized}
\dfrac{\lvert \varphi(\lambda)-\omega\rvert^{2}}{1-\lvert \varphi(\lambda) \rvert^{2}} \leq \alpha \max_{\lvert \tau_{j} \rvert=1} \dfrac{\lvert \lambda_{j}-\tau_{j}\rvert^{2}}{1-\lvert \lambda_{j} \rvert^{2}}.
\end{equation}
\end{theorem}

 \begin{align}\label{phi-omega}
 \lvert \varphi(\lambda)-\omega\rvert^{2}& \leq \|x\|_{\mathcal{M}}^{2}\sum_{j=1}^{d}\lvert 1-\overline{\tau}_{j}\lambda_{j}\rvert^{2}\|u_{j}(\lambda)\|_{\mathcal{M}_{j}}^{2}\nonumber\\
&= \|x\|_{\mathcal{M}}^{2}\sum_{j=1}^{d}\lvert \tau_{j}-\lambda_{j}\rvert^{2}\dfrac{1-\lvert \lambda_{j}\rvert^{2}}{1-\lvert \lambda_{j}\rvert^{2}}\|u_{j}(\lambda)\|_{\mathcal{M}_{j}}^{2}\nonumber\\
&\leq \|x\|_{\mathcal{M}}^{2}\bigg(\max_{j=1,\hdots,d}\dfrac{\lvert \tau_{j}-\lambda_{j}\rvert^{2}}{1-\lvert \lambda_{j}\rvert^{2}}\bigg)\sum_{j=1}^{d}\Big({1-\lvert \lambda_{j}\rvert^{2}}\Big)\|u_{j}(\lambda)\|_{\mathcal{M}_{j}}^{2}\nonumber\\
&=\|x\|_{\mathcal{M}}^{2}\bigg(\max_{j=1,\hdots,d}\dfrac{\lvert \tau_{j}-\lambda_{j}\rvert^{2}}{1-\lvert \lambda_{j}\rvert^{2}}\bigg)(1-\lvert \varphi(\lambda) \rvert^{2}),
\end{align}

\begin{proposition}\label{phi-nt-omega}
    If $\varphi\in\mathcal{SA}_{d}$ and $\tau\in\partial\mathbb{D}^{d}$ is a $B$-point for $\varphi$ then there exists $\omega\in\mathbb{T}$ such that 
 $$ \lim_{\lambda{\overset{nt}\rightarrow} \tau} \varphi(\lambda)= \omega.$$  
\end{proposition}

\begin{corollary}\label{coro5.7cara}
Let $\varphi\in\mathcal{SA}_{d}$, let $\tau\in\partial\mathbb{D}^{d}$ and let $(\mathcal{M},u)$ be a model of $\varphi$. \\
{\rm (i)} If $ \ \tau \in \partial \D^d \setminus\T^d$, then  the following conditions are equivalent: 
\begin{enumerate}
    \item $\dfrac{1-\lvert \varphi(\lambda) \rvert}{1-\|\lambda\|_{\infty}}$ is bounded on some sequence $\{\lambda^{(n)}\}$ that converges to $\tau$;
    \item $\lim_{n \to \infty} |\varphi(\lambda^{(n)})| =1$ and $u(\lambda^{(n)})$ is bounded on some sequence $\{\lambda^{(n)}\}$ that approaches $\tau$ nontangentially;
    \item $\dfrac{1-\lvert \varphi(\lambda) \rvert}{1-\|\lambda\|_{\infty}}$ is bounded on every subset of $\mathbb{D}^{d}$ that approaches $\tau$ nontangentially;
    \item for every subset $S$ of $\D^d$ that approaches $\tau$ nontangentially,  $\lim_{\la\to\tau, \la\in S} |\varphi(\lambda)| =1$, and  $ u(\lambda) $ is  bounded  on $S$.
\end{enumerate} 
{\rm (ii)} If $ \ \tau \in \T^d$, then  the following conditions are equivalent: 
\begin{enumerate}
    \item $\dfrac{1-\lvert \varphi(\lambda) \rvert}{1-\|\lambda\|_{\infty}}$ is bounded on some sequence $\{\lambda^{(n)}\}$ that converges to $\tau$;
    \item $u(\lambda^{(n)})$ is bounded on some sequence $\{\lambda^{(n)}\}$ that approaches $\tau$ nontangentially;
    \item $\dfrac{1-\lvert \varphi(\lambda) \rvert}{1-\|\lambda\|_{\infty}}$ is bounded on every subset of $\mathbb{D}^{d}$ that approaches $\tau$ nontangentially;
    \item $u(\lambda)$ is bounded on every subset of $\mathbb{D}^{d}$ that approaches $\tau$ nontangentially.
\end{enumerate}
\end{corollary}

\begin{lemma}\label{ujnotxj}
Let $\varphi\in\mathcal{SA}_{d}$, let $(\mathcal{M},u)$ be a model of $\varphi$, where  $\mathcal{M}$ is a separable complex Hilbert space with an orthogonal decomposition $\mathcal{M}=\mathcal{M}_{1}\oplus\hdots\oplus\mathcal{M}_{d}$ and $u=(u_1, \dots, u_d)$ is an analytic map  $u_j: \mathbb{D}^{d} \to \mathcal{M}_{j} $ for $j =1, \dots, d$, such that condition 
\eqref{modeleqn}    holds, and let $\tau\in\partial\mathbb{D}^{d}$ be a $B$-point for $\varphi$. Suppose that $\{\lambda^{(n)}\}$ converges to $\tau$ nontangentially in $\mathbb{D}^{d}$ and $\{u(\lambda^{(n)})\}$ converges weakly to $x$ in $\mathcal{M}$ as $n\rightarrow \infty$.
Suppose also that, for some index $j$,  $\{u_j(\lambda^{(n)})\}$ does not tend to $x_j$  in norm as $n\rightarrow \infty$.
Then there exists a   subsequence $(\lambda^{(n_k)})$ of $(\lambda^{(n)})$ such that, 
\begin{equation} \label{lemma-uj-0}  \|u_{j}(\lambda^{(n_k)})\|_{\mathcal{M}_{j}}^{2}-\mathrm{Re} \langle u_{j}(\lambda^{(n_k)}),x_{j}\rangle_{\mathcal{M}_{j}} \geq 0 \;\ \text{for all} \; \ k \ge 1.
\end{equation}
\end{lemma}

\begin{proposition}\label{prop5.8cara}
Let $\varphi\in\mathcal{SA}_{d}$, let $(\mathcal{M},u)$ be a model of $\varphi$,  where  $\mathcal{M}$ is a separable complex Hilbert space with an orthogonal decomposition $\mathcal{M}=\mathcal{M}_{1}\oplus\hdots\oplus\mathcal{M}_{d}$, 
 and let $\tau\in\partial\mathbb{D}^{d}$ be a $B$-point for $\varphi$. Suppose that $\{\lambda^{(n)}\}$ converges to $\tau$ nontangentially in $\mathbb{D}^{d}$. If $\{u(\lambda^{(n)})\}$ converges weakly in $\mathcal{M}$ as $n\rightarrow \infty$, then $\{u(\lambda^{(n)})\}$ converges in norm as $n\rightarrow \infty$.
\end{proposition}

\begin{proof} Suppose  $u(\lambda^{(n)})\rightarrow x$ weakly in $\mathcal{M}$.  The equality
$$\|u(\lambda^{(n)})-x\|_{\mathcal{M}}^{2} = \langle u(\lambda^{(n)})-x,u(\lambda^{(n)})-x\rangle_{\mathcal{M}}$$
implies that
\begin{align}
\|u(\lambda^{(n)})-x\|_{\mathcal{M}}^{2}&=\|u(\lambda^{(n)})\|_{\mathcal{M}}^{2}-2\mathrm{Re} \langle u(\lambda^{(n)}),x\rangle_{\mathcal{M}} +\|x\|_{\mathcal{M}}^{2}\nonumber\\
&=\bigg(\|u(\lambda^{(n)})\|_{\mathcal{M}}^{2}-\mathrm{Re} \langle u(\lambda^{(n)}),x\rangle_{\mathcal{M}} \bigg)\nonumber\\
&\hspace{3cm}+\bigg(\|x\|_{\mathcal{M}}^{2}-\mathrm{Re} \langle u(\lambda^{(n)}),x\rangle_{\mathcal{M}} \bigg)\label{2ndequationinthisalign}.
\end{align}
Since $u(\lambda^{(n)})\rightarrow x$ weakly as $n\rightarrow\infty$, 
$$\langle u(\lambda^{(n)}),x \rangle_{\mathcal{M}} \rightarrow \langle x,x\rangle_{\mathcal{M}}~~\text{as}~n\rightarrow \infty.$$
Thus
\begin{equation}\label{equivalenceofweakconvergence}
 \|x\|_{\mathcal{M}}^{2}-\mathrm{Re} \langle u(\lambda^{(n)}),x\rangle_{\mathcal{M}} \rightarrow 0~\text{as}~n\rightarrow \infty.
\end{equation}
To show that $u(\lambda^{(n)})$ converges in norm to $x$ as $n\rightarrow \infty$, all that remains is to show that the first term on the right hand side of equation (\ref{2ndequationinthisalign}), that is, 
\begin{equation}\label{ujtoxj}
\sum_{j=1}^{d}\Big(\|u_{j}(\lambda^{(n)})\|_{\mathcal{M}_{j}}^{2}-\mathrm{Re} \langle u_{j}(\lambda^{(n)}),x_{j}\rangle \Big) \to 0
\end{equation}
as  $n\rightarrow \infty$.
Since $\{u(\lambda^{(n)})\} \to x$ weakly in $\mathcal{M}$ as $n\rightarrow \infty$, it follows that, for each index $j= 1,2, \dots, d$,  $\{u_{j}(\lambda^{(n)})\} \to x_{j}$ weakly in $\mathcal{M}_{j}$ as $n\rightarrow \infty$.
 If for each $j$, $\|u_{j}(\lambda^{(n)}) - x_j\| \to 0$  as $n\rightarrow \infty$, then $u(\lambda^{(n)})$ converges in norm to $x$ as $n\rightarrow \infty$. Suppose that there are some $j$ such that 
$\|u_{j}(\lambda^{(n)}) - x_j\|$ does not tend to $0$  as $n\rightarrow \infty$. Let us show that in this case we will get a contradiction.

 By Proposition \ref{phi-nt-omega},
there exists $\omega\in\mathbb{T}$ such that $\varphi(\lambda)\rightarrow \omega~~\text{as}~~\lambda{\overset{nt}\rightarrow}\tau.$ By Proposition \ref{prop4.2generalized} (3), 
\begin{equation}\label{omega-model}
  1-\overline{\omega}\varphi(\lambda)=\sum_{\lvert \tau_{j} \rvert = 1}\Big(1-\overline{\tau}_{j}\lambda_{j}\Big)\langle u_{j}(\lambda),x_{j}\rangle_{\mathcal{M}_{j}}
  \end{equation}
holds for all $\lambda\in\mathbb{D}^{d}$.
By Proposition \ref{prop4.2generalized} (2), if $(x_{1},\hdots,x_{d})=x\in Y_{\tau}$ where $\tau\in\partial\mathbb{D}^{d}$ is such that  $(\tau_{1},\hdots,\tau_{d})=\tau\in\mathbb{T}^{a}\times\mathbb{D}^{b}$ where $a+b=d$, then for each $\tau_{j}\in\mathbb{D}$, $x_{j}=0$.

Thus,  by equation \eqref{omega-model},
   for $\lambda^{(n)}\in\mathbb{D}^{d}$, $n=1,2,\dots$,
\begin{equation}\label{prop4.2equationhere}
1-\overline{\omega}\varphi(\lambda^{(n)})=\sum_{j=1}^{d}\Big(1-\overline{\tau}_{j}\lambda^{(n)}_{j}\Big)\langle u_{j}(\lambda^{(n)}),x_{j}\rangle_{\mathcal{M}_{j}}.
\end{equation}
Hence,
\begin{equation}\label{REprop4.2equationhere}
\mathrm{Re}\bigg(1-\overline{\omega}\varphi(\lambda^{(n)})\bigg)=\mathrm{Re}\bigg(\sum_{j=1}^{d}\Big(1-\overline{\tau}_{j}\lambda^{(n)}_{j}\Big)\langle u_{j}(\lambda^{(n)}),x_{j}\rangle_{\mathcal{M}_{j}}\bigg).
\end{equation}
By equations (\ref{whenlambdametmu}) and (\ref{REprop4.2equationhere}),
\begin{align}
&\sum_{j=1}^{d}(1-\lvert \lambda_{j}^{(n)} \rvert^{2})\|u_{j}(\lambda^{(n)})\|_{\mathcal{M}_{j}}^{2}\nonumber\\
&= 1-\lvert \varphi(\lambda^{(n)}) \rvert^{2} \nonumber\\
&= 1-\lvert \varphi(\lambda^{(n)}) \rvert^{2}-2\mathrm{Re}\big(1-\overline{\omega}\varphi(\lambda^{(n)})\big)\nonumber\\
&\hspace{4cm}+ \mathrm{Re}\bigg(\sum_{j=1}^{d}2\Big(1-\overline{\tau}_{j}\lambda^{(n)}_{j}\Big)\langle u_{j}(\lambda^{(n)}),x_{j}\rangle_{\mathcal{M}_{j}}\bigg).
\label{rearrangmentofthemodelhere}
\end{align}
Subtraction of  $\; \sum_{j=1}^{d}(1-\lvert\lambda^{(n)}_{j}\rvert^{2})\mathrm{Re} \langle u_{j}(\lambda^{(n)}),x_{j}\rangle_{\mathcal{M}_{j}} \;$ from both sides of equation (\ref{rearrangmentofthemodelhere}) gives us 
\begin{align}
&\sum_{j=1}^{d}(1-\lvert \lambda^{(n)}_{j} \rvert^{2})\Big(\|u_{j}(\lambda^{(n)})\|_{\mathcal{M}_{j}}^{2}-\mathrm{Re}\langle u_{j}(\lambda^{(n)}),x_{j}\rangle  \Big)\nonumber\\
&=1-\lvert \varphi(\lambda^{(n)}) \rvert^{2}-2\mathrm{Re}\big(1-\overline{\omega}\varphi(\lambda^{(n)})\big)\nonumber\\
&\hspace{2cm}+ \mathrm{Re}\bigg(\sum_{j=1}^{d}2\Big(1-\overline{\tau}_{j}\lambda^{(n)}_{j}\Big)\langle u_{j}(\lambda^{(n)}),x_{j}\rangle_{\mathcal{M}_{j}}\bigg) - \sum_{j=1}^{d}(1-\lvert\lambda^{(n)}_{j}\rvert^{2})
\mathrm{Re} \langle u_{j}(\lambda^{(n)}),x_{j}\rangle_{\mathcal{M}_{j}} \nonumber \\
&=-1+2\mathrm{Re}\big(\overline{\omega}\varphi(\lambda^{(n)})\big)-\lvert \varphi(\lambda^{(n)}) \rvert^{2}\nonumber \\
&\hspace{2cm}+ \mathrm{Re}\bigg[\sum_{j=1}^{d}\Big(2(1-\overline{\tau}_{j}\lambda^{(n)}_{j})-(1-\lvert \lambda^{(n)}_{j}\rvert^{2})\Big)\langle u_{j}(\lambda^{(n)}),x_{j}\rangle_{\mathcal{M}_{j}}\bigg]. \label{Prop4.5-01}
\end{align} 
Note that
$$\lvert1-\overline{\omega}\varphi(\lambda^{(n)})\rvert^{2}= 1-2\mathrm{Re}\big(\overline{\omega}\varphi(\lambda^{(n)})\big)+\lvert \varphi(\lambda^{(n)}) \rvert^{2}, $$
and so by equation \eqref{Prop4.5-01},
\begin{align}
&\sum_{j=1}^{d}(1-\lvert \lambda^{(n)}_{j} \rvert^{2})\Big(\|u_{j}(\lambda^{(n)})\|_{\mathcal{M}_{j}}^{2}-
\mathrm{Re} \langle u_{j}(\lambda^{(n)}),x_{j}\rangle \Big)\nonumber\\
&=-\lvert1-\overline{\omega}\varphi(\lambda^{(n)})\rvert^{2}+\mathrm{Re}\Big[\sum_{j=1}^{d}\Big(1-2\overline{\tau}_{j}\lambda^{(n)}_{j}+\lvert\lambda^{(n)}_{j}\rvert^{2}\Big)\|x_{j}\|_{\mathcal{M}_{j}}^{2}\Big]\nonumber\\
&\hspace{2.5cm}+\mathrm{Re}\bigg[\sum_{j=1}^{d}\Big(1-2\overline{\tau}_{j}\lambda^{(n)}_{j}+\lvert\lambda^{(n)}_{j}\rvert^{2}\Big)\langle u_{j}(\lambda^{(n)})-x_{j},x_{j}\rangle_{\mathcal{M}_{j}}\bigg].\label{lastequationhereinthispart}
\end{align} 
Note that, for $\lambda\in\mathbb{D}^{d}$, $\lvert \lambda_{j}\rvert^{2}\leq \lvert \lambda_{j} \rvert \leq \|\lambda\|_{\infty} <1$ for $j=1,\hdots, d$, and so
$$ 0< \dfrac{1}{1-\lvert \lambda_{j}\rvert^{2}} \leq \dfrac{1}{1-\|\lambda\|_{\infty}} < \infty.$$
By assumption, $\lambda^{(n)}{\overset{nt}\rightarrow}\tau$ as $n\rightarrow \infty$. Hence there exists a $c>0$ such that
$$\|\lambda^{(n)} - \tau\|_{\infty}\leq c(1-\|\lambda^{(n)}\|_{\infty}) \;\; \text{for all } \; n \in \mathbb{N}.$$
Consequently, for all $ n \in \mathbb{N}$,
$$\dfrac{\|\lambda^{(n)}-\tau\|_{\infty}}{1-\lvert \lambda^{(n)}_{j} \rvert^{2}}\leq \dfrac{\|\lambda^{(n)}-\tau\|_{\infty}}{1-\|\lambda^{(n)}\|_{\infty}}\leq c,$$
so that
\begin{equation}\label{inequalitywithc}
\dfrac{c \big(1-\lvert \lambda^{(n)}_{j}\rvert^{2}\big)}{\|\lambda^{(n)}-\tau\|_{\infty}}\geq 1.
\end{equation}
It is easy to see that, for all $\lambda\in\mathbb{D}^{d}$ and $\tau \in\partial\mathbb{D}^{d}$,
$$1-\|\lambda\|_{\infty}\leq \|\lambda-\tau\|_{\infty},$$
and so
\begin{equation}\label{inequalitywithc-5}
\dfrac{1}{\|\lambda-\tau\|_{\infty}} \leq \dfrac{1}{1-\|\lambda\|_{\infty}}.
\end{equation}
Note that, for $j=1,\hdots, d$,
$$1-\lvert \lambda_{j}\rvert^{2} \leq \|\tau+\lambda\|_{\infty}\| \tau-\lambda\|_{\infty}.$$
Let $c_{1}=\max_{i=1,\hdots d} \lvert \tau_{i}+\lambda_{i}\rvert = \|\tau+\lambda\|_{\infty}$. Then, for $j=1,\hdots, d$, we have the inequality
\begin{equation}\label{ineqc1}
1-\lvert \lambda^{(n)}_{j}\rvert^{2} \leq c_{1}\|\tau -\lambda^{(n)}\|_{\infty}.
\end{equation}
Equivalently, for $j=1,\hdots, d$,
\begin{equation}\label{inequalityforc1}
\dfrac{1}{\|\tau-\lambda^{(n)}\|_{\infty}} \leq \dfrac{c_{1}}{1-\lvert \lambda^{(n)}_{j} \rvert^{2}}.
\end{equation}
Multiply equation \eqref{lastequationhereinthispart} by $\dfrac{c}{\|\lambda^{(n)}-\tau\|_{\infty}}$, to obtain the following equality, 
\begin{align}\label{equation527}
& \dfrac{c}{\|\lambda^{(n)}-\tau\|_{\infty}}\cdot\sum_{j=1}^{d}(1-\lvert \lambda^{(n)}_{j} \rvert^{2})\Big(\|u_{j}(\lambda^{(n)})\|_{\mathcal{M}_{j}}^{2}-\mathrm{Re} \langle u_{j}(\lambda^{(n)}),x_{j}\rangle_{\mathcal{M}_{j}} \Big)\nonumber\\
&=\dfrac{c}{\|\lambda^{(n)}-\tau\|_{\infty}} \Bigg(-\lvert1-\overline{\omega}\varphi(\lambda^{(n)})\rvert^{2}+\sum_{j=1}^{d}\mathrm{Re}\Big(1-2\overline{\tau}_{j}\lambda^{(n)}_{j}+\lvert\lambda^{(n)}_{j}\rvert^{2}\Big)\|x_{j}\|_{\mathcal{M}_{j}}^{2}
\nonumber\\
&\hspace{1cm}+\mathrm{Re}\bigg[\sum_{j=1}^{d}\Big(1-2\overline{\tau}_{j}\lambda^{(n)}_{j}+\lvert\lambda^{(n)}_{j}\rvert^{2}\Big)\langle u_{j}(\lambda^{(n)})-x_{j},x_{j}\rangle_{\mathcal{M}_{j}}\bigg]\Bigg)\nonumber\\
&=c\Bigg(\dfrac{-\lvert1-\overline{\omega}\varphi(\lambda^{(n)})\rvert^{2}}{\|\lambda^{(n)}-\tau\|_{\infty}}+\sum_{j=1}^{d}\mathrm{Re}\Big(1-2\overline{\tau}_{j}\lambda^{(n)}_{j}+\lvert\lambda^{(n)}_{j}\rvert^{2}\Big)\dfrac{\|x_{j}\|_{\mathcal{M}_{j}}^{2}}{\|\lambda^{(n)}-\tau\|_{\infty}}\nonumber\\
&\hspace{1cm}+\sum_{j=1}^{d}\dfrac{1}{\|\lambda^{(n)}-\tau\|_{\infty}}\mathrm{Re}\bigg[
\Big(1-2\overline{\tau}_{j}\lambda^{(n)}_{j}+\lvert\lambda^{(n)}_{j}\rvert^{2}\Big)\langle u_{j}(\lambda^{(n)})-x_{j},x_{j}\rangle_{\mathcal{M}_{j}}\bigg]\Bigg).
\end{align}

There are two cases: (i) $\tau\in\mathbb{T}^{d}$ and (ii) $\tau\in\partial\mathbb{D}^{d}\setminus\mathbb{T}^{d}$.

{\rm (i)}. Assume firstly that $\tau\in\mathbb{T}^{d}$. Then, for each $j$, 
\begin{equation}\label{inequalityfor527}
\mathrm{Re}\Big( 1-2\overline{\tau}_{j}\lambda^{(n)}_{j}+\lvert\lambda^{(n)}_{j}\rvert^{2} \Big)= |{\tau}_{j} -\lambda^{(n)}_{j}|^2,
\end{equation}
and 
\begin{align}\label{inequalityfor527-2}
&\Big|1-2\overline{\tau}_{j}\lambda^{(n)}_{j}+\lvert\lambda^{(n)}_{j}\rvert^{2}\Big| = \Big| \overline{\tau}_{j} {\tau}_{j}- \overline{{\tau}_{j}}\lambda^{(n)}_{j} -\overline{{\tau}_{j}}\lambda^{(n)}_{j} + \overline{\lambda^{(n)}_{j}}\lambda^{(n)}_{j}\Big|\nonumber\\
&\leq \Big|{\tau}_{j} -\lambda^{(n)}_{j}\Big| + \Big|\lambda_{j}^{(n)} \Big|\Big|\overline{{\tau}_{j}} -\overline{\lambda^{(n)}_{j}}\Big| = \Big|{\tau}_{j} -\lambda^{(n)}_{j}\Big| \Big(1 + \Big|\lambda_{j}^{(n)} \Big| \Big).
\end{align} 
Thus, by equations  (\ref{equation527}), \eqref{inequalityfor527} and inequalities \eqref{inequalityfor527-2} and \eqref{inequalitywithc-5},
\begin{align}\label{equation5271}
& \Bigg\lvert \dfrac{c}{\|\lambda^{(n)}-\tau\|_{\infty}}\cdot\sum_{j=1}^{d}(1-\lvert \lambda^{(n)}_{j} \rvert^{2})\Big(\|u_{j}(\lambda^{(n)})\|_{\mathcal{M}_{j}}^{2}-\mathrm{Re} \langle u_{j}(\lambda^{(n)}),x_{j}\rangle_{\mathcal{M}_{j}} \Big)
\Bigg\rvert \nonumber\\
&\leq c\Bigg(\dfrac{\lvert 1-\overline{\omega}\varphi(\lambda^{(n)})\rvert^{2}}{\|\lambda^{(n)}-\tau\|_{\infty}}+
\sum_{j=1}^{d}\dfrac{\Big[|{\tau}_{j} -\lambda^{(n)}_{j}|^2 
\|x_{j}\|_{\mathcal{M}_{j}}^{2}\Big]}{\|\lambda^{(n)}-\tau\|_{\infty}}\nonumber\\
&\hspace{4cm}+\sum_{j=1}^{d}\dfrac{\Big[
{\lvert \tau_{j}-\lambda^{(n)}_{j}\rvert}{(1+\lvert\lambda^{(n)}_{j})}
|\langle u_{j}(\lambda^{(n)})-x_{j},x_{j}\rangle_{\mathcal{M}_{j}}|\Big]}{\|\lambda^{(n)}-\tau\|_{\infty}}\nonumber\\
&\leq c\Bigg(\dfrac{\lvert 1-\overline{\omega}\varphi(\lambda^{(n)})\rvert^{2}}{1-\|\lambda^{(n)}\|_{\infty}}
+\|x\|_{\mathcal{M}}^{2}\sum_{j=1}^{d}\dfrac{\Big[|{\tau}_{j} -\lambda^{(n)}_{j}|^2 \Big]}{\|\lambda^{(n)}-\tau\|_{\infty}} 
\nonumber\\
&\hspace{4cm}+\sum_{j=1}^{d}\dfrac{\Big[
{\lvert \tau_{j}-\lambda^{(n)}_{j}\rvert}{(1+\lvert\lambda^{(n)}_{j})}
\Big]}{\|\lambda^{(n)}-\tau\|_{\infty}}\lvert\langle u_{j}(\lambda^{(n)})-x_{j},x_{j}\rangle_{\mathcal{M}}\rvert \Bigg).
\end{align} 
The first term on the right tends to zero because, firstly, by inequality \eqref{phi-omega} from Theorem \ref{theorem4.9generalized},
\begin{equation}\label{phi-omega-2}
 \lvert \varphi(\lambda)-\omega\rvert^{2} \leq  \|x\|_{\mathcal{M}}^{2}\bigg(\max_{j=1,\hdots,d}\dfrac{\lvert \tau_{j}-\lambda_{j}\rvert^{2}}{1-\lvert \lambda_{j}\rvert^{2}}\bigg)(1-\lvert \varphi(\lambda) \rvert^{2}),
\end{equation}
 secondly, by Corollary \ref{coro5.7cara},   $\dfrac{1-\lvert \varphi(\lambda) \rvert}{1-\|\lambda\|_{\infty}}$ is bounded on every subset of $\mathbb{D}^{d}$ that approaches $\tau$ nontangentially and because $\lambda^{(n)}{\overset{nt}\rightarrow}\tau$. 
The second term tends to zero because $\lambda^{(n)}{\overset{nt}\rightarrow}\tau$. The third term tends to zero because
each summand is the product of a bounded factor and a factor that tends to zero, since $u(\lambda^{(n)})\rightarrow x$ weakly in $\mathcal{M}$. Therefore,
\begin{equation}\label{equation5272}
 \dfrac{c}{\|\lambda^{(n)}-\tau\|_{\infty}}\cdot\sum_{j=1}^{d}(1-\lvert \lambda^{(n)}_{j} \rvert^{2})\Big(\|u_{j}(\lambda^{(n)})\|_{\mathcal{M}_{j}}^{2}-\mathrm{Re} \langle u_{j}(\lambda^{(n)}),x_{j}\rangle_{\mathcal{M}_{j}} \Big) \to 0
\end{equation}
as $n \rightarrow  \infty$.

Let $$\Lambda = \{ j: \ \text{such that} \ \|u_{j}(\lambda^{(n)}) - x_j\| \; \text{ does not tend to} \; 0 \; \text{as} \; n\rightarrow \infty \},$$
and let $$\Lambda'= \{ j: \ \text{such that} \ \|u_{j}(\lambda^{(n)}) - x_j\| \; \text{tends to} \; 0 \; \text{as} \; n \rightarrow \infty \}.$$
 By Lemma \ref{ujnotxj} and by repeatedly passing to subsequences, we can choose a strictly increasing sequence 
 $(n_k)_{k \ge 1}$ such that,  for each $j \in \Lambda$, 
$$\Big(\|u_{j}(\lambda^{(n_k)})\|_{\mathcal{M}_{j}}^{2}-\mathrm{Re} \langle u_{j}(\lambda^{(n_k)}),x_{j}\rangle_{\mathcal{M}_{j}} \Big) \geq 0  \; \text{ for all } \; k \ge1.$$
By inequality (\ref{inequalitywithc}), 
\begin{align}
& \sum_{j \in \Lambda}\Big(\|u_{j}(\lambda^{(n_k)})\|_{\mathcal{M}_{j}}^{2}-\mathrm{Re} \langle u_{j}(\lambda^{(n_k)}),x_{j}\rangle_{\mathcal{M}_{j}} \Big) \nonumber\\
& \leq \sum_{j \in \Lambda} \dfrac{c (1-\lvert \lambda^{(n_k)}_{j} \rvert^{2})}{\|\lambda^{(n_k)}-\tau\|_{\infty}} 
\Big(\|u_{j}(\lambda^{(n_k)})\|_{\mathcal{M}_{j}}^{2}-\mathrm{Re} \langle u_{j}(\lambda^{(n_k)}),x_{j}\rangle_{\mathcal{M}_{j}} \Big)\nonumber\\
&= \dfrac{c}{\|\lambda^{(n_k)}-\tau\|_{\infty}}\sum_{j \in \Lambda}(1-\lvert \lambda^{(n_k)}_{j} \rvert^{2})\Big(\|u_{j}(\lambda^{(n_k)})\|_{\mathcal{M}_{j}}^{2}-\mathrm{Re} \langle u_{j}(\lambda^{(n_k)}),x_{j}\rangle_{\mathcal{M}_{j}} \Big).\label{oneovercinequality}
\end{align}
Therefore, by inequality \eqref{oneovercinequality} and relation \eqref{equation5272},
\begin{align}
& \sum_{j \in \Lambda}\Big(\|u_{j}(\lambda^{(n_k)})\|_{\mathcal{M}_{j}}^{2}-\mathrm{Re} \langle u_{j}(\lambda^{(n_k)}),x_{j}\rangle_{\mathcal{M}_{j}} \Big) \nonumber\\
& \hspace{1cm} +\dfrac{c}{\|\lambda^{(n_k)}-\tau\|_{\infty}}\sum_{j \in \Lambda'}(1-\lvert \lambda^{(n_k)}_{j} \rvert^{2})\Big(\|u_{j}(\lambda^{(n_k)})\|_{\mathcal{M}_{j}}^{2}-\mathrm{Re} \langle u_{j}(\lambda^{(n_k)}),x_{j}\rangle_{\mathcal{M}_{j}} \Big)\nonumber\\
&\leq \dfrac{c}{\|\lambda^{(n_k)}-\tau\|_{\infty}}\cdot\sum_{j=1}^{d}(1-\lvert \lambda^{(n_k)}_{j} \rvert^{2})\Big(\|u_{j}(\lambda^{(n_k)})\|_{\mathcal{M}_{j}}^{2}-\mathrm{Re} \langle u_{j}(\lambda^{(n_k)}),x_{j}\rangle_{\mathcal{M}_{j}} \Big)
 \to 0  \label{oneovercinequality2}
\end{align}
 as $ n_k \rightarrow  \infty$.
Since, for each $j \in \Lambda'$, $\|u_{j}(\lambda^{(n)}) - x_j\| \rightarrow 0 $ as $n \rightarrow \infty,$
by inequality \eqref{ineqc1},
$$\dfrac{c}{\|\lambda^{(n_k)}-\tau\|_{\infty}}\sum_{j \in \Lambda'}(1-\lvert \lambda^{(n_k)}_{j} \rvert^{2})\Big(\|u_{j}(\lambda^{(n_k)})\|_{\mathcal{M}_{j}}^{2}-\mathrm{Re} \langle u_{j}(\lambda^{(n_k)}),x_{j}\rangle_{\mathcal{M}_{j}} \Big)
\rightarrow 0$$
 as  $n_k\rightarrow \infty$. Hence
$$ \sum_{j \in \Lambda}\Big(\|u_{j}(\lambda^{(n_k)})\|_{\mathcal{M}_{j}}^{2}-\mathrm{Re} \langle u_{j}(\lambda^{(n_k)}),x_{j}\rangle_{\mathcal{M}_{j}} \Big)\rightarrow 0 \; \text {as} \; n_k \rightarrow \infty,$$
contrary to our assumption. Therefore, for each $j$, 
$$\|u_{j}(\lambda^{(n)}) - x_j\| \rightarrow 0 \; \text{as} \; n \rightarrow \infty,$$
and so
 $$\|u(\lambda^{(n)})\|_{\mathcal{M}}^{2}-\mathrm{Re} \langle u(\lambda^{(n)}),x\rangle_{\mathcal{M}} \rightarrow 0 \; \text {as} \; n \rightarrow \infty.$$  
Therefore,  $u(\lambda^{(n)})$ converges in norm to $x$ as $n \rightarrow \infty$ when $\tau\in\mathbb{T}^{d}$.
 

 {\rm (ii)}. For the second case, let $a+b=d$ such that $\tau\in\mathbb{T}^{a}\times \mathbb{D}^{b}$. By Proposition \ref{prop4.2generalized}, we have for $j=b,\hdots, d$, $x_{j}=0$. Indeed, as 
$$1-\lvert \varphi(\lambda^{(n)})\rvert^{2}=\sum_{j=1}^{d}\big(1-\lvert \lambda_{j}^{(n)}\rvert^{2}\big)\|u_{j}(\lambda^{(n)})\|_{\mathcal{M}_{j}}^{2},$$
letting $\lambda^{(n)}$ tend to $\tau$ we get $1-\lvert \varphi(\lambda^{(n)})\rvert^{2}$ tends to $0$ and 
$\big(1-\lvert \lambda_{j}^{(n)}\rvert^{2}\big)$ tends to $1 -|\tau_j|^2 \neq 0$ for  all $j=b,\hdots, d$. Hence
 $\|u_{j}(\lambda^{(n)})\|_{\mathcal{M}_{j}}$ tends to zero for all $j=b,\hdots, d$. 
It is well known that in a Hilbert space, a sequence converges in norm if and only if it converges weakly and the norms converge in $\mathbb{R}$. Hence, for each $j$ such that $ |\tau_j| < 1$, $u_{j}(\lambda^{(n)})$ tends to  $x_{j}=0$ in norm in ${\mathcal{M}_{j}}$.

 Suppose again that there are some $j$ such that 
$\|u_{j}(\lambda^{(n)}) - x_j\|$ does not tend to $0$  as $n\rightarrow \infty$. Let us show that in this case we will also get a contradiction.
Let $$\Lambda = \{ j: \ \text{such that} \ \|u_{j}(\lambda^{(n)}) - x_j\| \; \text{ does not tend to} \; 0 \; \text{as} \; n\rightarrow \infty \},$$
and let $$\Lambda'= \{ j: \ \text{such that} \ \|u_{j}(\lambda^{(n)}) - x_j\| \; \text{tends to} \; 0 \; \text{as} \; n \rightarrow \infty \}.$$
We have shown that $\{ j : |\tau_j| < 1 \}$ is a subset of $\Lambda'$. By equality \eqref{equation527},
since $x_j=0$ for each $j$ such that $|\tau_j| < 1$, we have the following relation:
\begin{align}\label{equation527notTd}
& \dfrac{c}{\|\lambda^{(n)}-\tau\|_{\infty}}\cdot\sum_{j=1}^{d}(1-\lvert \lambda^{(n)}_{j} \rvert^{2})\Big(\|u_{j}(\lambda^{(n)})\|_{\mathcal{M}_{j}}^{2}-\mathrm{Re} \langle u_{j}(\lambda^{(n)}),x_{j}\rangle_{\mathcal{M}_{j}} \Big)\nonumber\\
&=c\Bigg(\dfrac{-\lvert1-\overline{\omega}\varphi(\lambda^{(n)})\rvert^{2}}{\|\lambda^{(n)}-\tau\|_{\infty}}+\sum_{j: |\tau_j| =1}\mathrm{Re}\Big(1-2\overline{\tau}_{j}\lambda^{(n)}_{j}+\lvert\lambda^{(n)}_{j}\rvert^{2}\Big)\dfrac{\|x_{j}\|_{\mathcal{M}_{j}}^{2}}{\|\lambda^{(n)}-\tau\|_{\infty}}\nonumber\\
&\hspace{1cm}+\sum_{j: |\tau_j| =1}\dfrac{1}{\|\lambda^{(n)}-\tau\|_{\infty}}\mathrm{Re}\bigg[
\Big(1-2\overline{\tau}_{j}\lambda^{(n)}_{j}+\lvert\lambda^{(n)}_{j}\rvert^{2}\Big)\langle u_{j}(\lambda^{(n)})-x_{j},x_{j}\rangle_{\mathcal{M}_{j}}\bigg]\Bigg).
\end{align}
If we simply repeat the calculations as before in  case (i), we produce the same result, that is, for $\tau\in\mathbb{T}^{a}\times\mathbb{D}^{b}$, for each $j$,
$$\|u_{j}(\lambda^{(n)})- x_{j}\|_{\mathcal{M}_{j}}\rightarrow 0~~\text{as}~~\lambda^{(n)}{\overset{nt}\rightarrow}\tau.$$
Therefore, for $\tau\in\partial\mathbb{D}^{d}$, $u(\lambda^{(n)})$ converges in norm to $x$ as $n \rightarrow \infty$.
\end{proof}

\end{document}
